\documentclass[11pt,a4paper]{article}

% --- Encoding ---
\usepackage[utf8]{inputenc}
\usepackage[T1]{fontenc}
\usepackage[spanish,es-noshorthands]{babel}
\usepackage{lmodern}

% --- Layout ---
\usepackage[margin=2.5cm]{geometry}
\usepackage{parskip}
\setlength{\parindent}{0pt}
\setlength{\parskip}{0.5em plus 0.2em minus 0.1em}

% --- Tables and graphics ---
\usepackage{booktabs}
\usepackage{array}
\usepackage{enumitem}
\usepackage{xcolor}
\usepackage{graphicx}
\usepackage{tikz}
\usetikzlibrary{positioning,arrows.meta,shapes,fit,backgrounds,calc}
\usepackage{caption}
\usepackage{subcaption}
\usepackage{amsmath,amssymb}

% --- Code listings ---
\usepackage{listings}
\definecolor{codegray}{rgb}{0.96,0.96,0.96}
\definecolor{codeborder}{rgb}{0.85,0.85,0.85}
\definecolor{codecomment}{rgb}{0.30,0.55,0.30}
\definecolor{codekeyword}{rgb}{0.10,0.30,0.70}
\definecolor{codestring}{rgb}{0.65,0.15,0.10}

\lstdefinestyle{shellstyle}{
    language=bash,
    backgroundcolor=\color{codegray},
    basicstyle=\ttfamily\small,
    breaklines=true,
    frame=single,
    rulecolor=\color{codeborder},
    showstringspaces=false,
    keywordstyle=\color{codekeyword}\bfseries,
    commentstyle=\color{codecomment}\itshape,
    stringstyle=\color{codestring},
    aboveskip=1em,
    belowskip=1em,
    columns=fullflexible,
}

\lstdefinestyle{textstyle}{
    backgroundcolor=\color{codegray},
    basicstyle=\ttfamily\footnotesize,
    breaklines=true,
    frame=single,
    rulecolor=\color{codeborder},
    showstringspaces=false,
    aboveskip=1em,
    belowskip=1em,
    columns=fullflexible,
}

\lstset{style=shellstyle}

% --- Hyperlinks ---
\usepackage[colorlinks=true,
            linkcolor=blue!60!black,
            urlcolor=blue!50!black,
            citecolor=blue!60!black]{hyperref}

% --- Gantt color palette ---
\definecolor{gp1}{HTML}{4C72B0}
\definecolor{gp2}{HTML}{DD8452}
\definecolor{gp3}{HTML}{55A868}
\definecolor{gp4}{HTML}{8172B2}

% --- Title block ---
\title{
    \vspace{-2.5em}
    \textbf{Práctico 4} \\[0.3em]
    \large Planificación de uso de CPU
}
\author{
    Tomás Rodeghiero \\
    \small Sistemas Operativos --- Universidad Nacional de Río Cuarto
}
\date{2026}

\begin{document}

\maketitle

\begin{abstract}
\noindent
Este documento presenta la resolución integral del Práctico 4 de la
materia Sistemas Operativos (UNRC, 2026), centrado en la
\emph{planificación de uso de CPU}. Se estudian seis ejercicios que
recorren los algoritmos clásicos: \emph{First-Come First-Served}
(FCFS), \emph{Shortest Job First} (SJF), \emph{Round Robin} (RR),
planificación por \emph{prioridades} estáticas y dinámicas y
\emph{Multilevel Feedback Queues} (MLF), más una discusión sobre
\emph{processor affinity} vs.\ \emph{load balancing} en sistemas SMP.
Cada ejercicio se justifica con la teoría del capítulo
\emph{Planificación de uso de CPU} de las Notas del curso (Marcelo
Arroyo) y se acompaña con tablas de procesos, diagramas de Gantt y
cálculos paso a paso de los tiempos de espera, turnaround y respuesta.
\end{abstract}

\tableofcontents
\bigskip

% =====================================================================
\section{Introducción}
% =====================================================================

El \emph{scheduler} de CPU es el subsistema del SO que decide
\emph{cuándo}, \emph{dónde} y \emph{por cuánto tiempo} cada tarea va
a ejecutar. La política que adopta depende del tipo de sistema y de
la carga de trabajo: en sistemas \emph{batch} se busca minimizar el
tiempo medio de espera; en \emph{time-sharing}, acotar el tiempo de
respuesta; en \emph{multiprocesadores SMP}, equilibrar la carga sin
romper la afinidad de caché. El capítulo de las Notas resume los
criterios habituales:

\begin{itemize}[leftmargin=*,nosep]
    \item Maximizar la utilización de la CPU.
    \item \textbf{Justicia (\emph{fairness})}: tiempos parejos para
        cada tarea.
    \item \textbf{Throughput}: procesos completados por unidad de
        tiempo.
    \item \textbf{Tiempo de retorno (\emph{turnaround})}: tiempo total
        que un proceso pasa en el sistema, desde que arribó hasta que
        terminó.
    \item \textbf{Tiempo de espera}: tiempo en estado \texttt{READY}
        sin ejecutar (CPU asignada a otro).
    \item \textbf{Tiempo de respuesta}: latencia desde que llega un
        requerimiento hasta que el sistema empieza a procesarlo.
\end{itemize}

Varios de estos criterios se contraponen entre sí: maximizar
\emph{throughput} puede deteriorar el tiempo de respuesta;
balancear la carga puede afectar la afinidad de caché. Por eso no
existe ``el'' mejor algoritmo, sino el más adecuado para cada
escenario.

Las fórmulas que se usan a lo largo del práctico son las estándar:

\begin{align*}
\text{Espera}_i      &= \text{Inicio}_i - \text{Arribo}_i
                       \quad (\text{para algoritmos no preemptivos sin I/O}) \\
\text{Turnaround}_i  &= \text{Fin}_i - \text{Arribo}_i \\
\text{Respuesta}_i   &= \text{PrimerInicio}_i - \text{Arribo}_i
\end{align*}

En la versión preemptiva con interrupciones por timer, el tiempo de
espera se calcula como
$\text{Espera}_i = \text{Turnaround}_i - \text{CPU}_i - \text{WAIT}_i$.

% =====================================================================
\section{Ejercicio 1 --- FCFS y SJF: tiempo de espera promedio}
% =====================================================================

\begin{center}
\small
\begin{tabular}{@{}lcc@{}}
\toprule
Proceso & Tiempo de creación (arribo) & Ráfaga de CPU (\(\mu\)s) \\
\midrule
P1 & 0.0 & 8 \\
P2 & 0.5 & 3 \\
P3 & 3.0 & 5 \\
\bottomrule
\end{tabular}
\end{center}

Se pide el tiempo de espera promedio bajo (1) FCFS con los tiempos de
arribo dados y (2) SJF \emph{asumiendo que los tres arriban en
$t = 0$}.

\subsection{1.a FCFS (FIFO) con los arribos dados}

FCFS es el algoritmo más simple: cola FIFO, no preemptivo. Se otorga
la CPU al primero de la cola y no se la quita hasta que termine. El
orden de arribo es P1, P2, P3.

\begin{center}
\begin{tikzpicture}[xscale=0.5,yscale=0.7]
    \draw[fill=gp1] (0,0) rectangle (8,1)  node[midway,white,font=\bfseries]{P1};
    \draw[fill=gp2] (8,0) rectangle (11,1) node[midway,white,font=\bfseries]{P2};
    \draw[fill=gp3] (11,0) rectangle (16,1) node[midway,white,font=\bfseries]{P3};
    \foreach \x in {0,8,11,16}
        \draw (\x,0) -- (\x,-0.15) node[below,font=\small]{\x};
\end{tikzpicture}
\\\small\emph{Diagrama de Gantt --- FCFS}
\end{center}

\begin{center}
\small
\begin{tabular}{@{}lcccc@{}}
\toprule
Proceso & Arribo & Inicio & Fin & Espera \\
\midrule
P1 & 0.0 & 0 & 8 & 0.0 \\
P2 & 0.5 & 8 & 11 & 7.5 \\
P3 & 3.0 & 11 & 16 & 8.0 \\
\bottomrule
\end{tabular}
\end{center}

\[
\overline{W}_{\text{FCFS}}
= \frac{0 + 7.5 + 8}{3} = \frac{15.5}{3} \approx 5.17.
\]

\textbf{Observación.} P1 es la ráfaga más larga y llega primero, así
que arrastra la espera de los demás (\emph{convoy effect}). Es
exactamente la desventaja que la teoría le adjudica a FCFS:
``comúnmente produce tiempos promedio de espera grandes''.

\subsection{1.b SJF (no preemptivo, todos en $t = 0$)}

Bajo el supuesto del enunciado, los tres procesos están disponibles
desde $t = 0$. SJF elige siempre la ráfaga más corta entre los
\texttt{READY}: P2 (3) primero, P3 (5) después, P1 (8) al final.

\begin{center}
\begin{tikzpicture}[xscale=0.5,yscale=0.7]
    \draw[fill=gp2] (0,0) rectangle (3,1) node[midway,white,font=\bfseries]{P2};
    \draw[fill=gp3] (3,0) rectangle (8,1) node[midway,white,font=\bfseries]{P3};
    \draw[fill=gp1] (8,0) rectangle (16,1) node[midway,white,font=\bfseries]{P1};
    \foreach \x in {0,3,8,16}
        \draw (\x,0) -- (\x,-0.15) node[below,font=\small]{\x};
\end{tikzpicture}
\\\small\emph{Diagrama de Gantt --- SJF}
\end{center}

\begin{center}
\small
\begin{tabular}{@{}lcccc@{}}
\toprule
Proceso & Arribo & Inicio & Fin & Espera \\
\midrule
P2 & 0 & 0 & 3  & 0 \\
P3 & 0 & 3 & 8  & 3 \\
P1 & 0 & 8 & 16 & 8 \\
\bottomrule
\end{tabular}
\end{center}

\[
\overline{W}_{\text{SJF}}
= \frac{0 + 3 + 8}{3} = \frac{11}{3} \approx 3.67.
\]

\subsection*{Comparación}

\begin{center}
\small
\begin{tabular}{@{}lll@{}}
\toprule
Algoritmo & Orden & Espera promedio \\
\midrule
FCFS & P1 - P2 - P3 & 5.17 \\
SJF (todos en $t=0$) & P2 - P3 - P1 & \textbf{3.67} \\
\bottomrule
\end{tabular}
\end{center}

SJF da menor tiempo de espera promedio. La razón es la propiedad
clásica del algoritmo: \emph{dado un conjunto fijo de ráfagas
conocidas, ejecutar primero las más cortas minimiza la espera
promedio}. Esta propiedad de optimalidad es la motivación histórica
de SJF en sistemas \emph{batch} con \emph{long-term schedulers},
donde el usuario declara el tiempo estimado del job.

% =====================================================================
\section{Ejercicio 2 --- Round Robin y comparación con FCFS y SJF}
% =====================================================================

Para los procesos del Ejercicio~1 se aplica \textbf{Round Robin}
preemptivo con quantum $q = 2$. Se mantiene una cola FIFO de
\texttt{RUNNABLE}: cuando un proceso llega o vuelve por consumir su
quantum se encola al final; el scheduler toma siempre el primero de
la cola.

\paragraph{Convención de empates.} Cuando un nuevo arribo y un
desalojo ocurren al mismo instante, primero se encola el recién
llegado y luego el desalojado. Es la convención habitual de los
textos clásicos (Silberschatz, Tanenbaum).

\subsection*{Traza paso a paso}

Notación: la cola se escribe \texttt{[cabeza ... cola]} y entre
paréntesis aparece la ráfaga remanente.

\begin{center}
\small
\begin{tabular}{@{}llp{6cm}l@{}}
\toprule
Intervalo & CPU & Evento al final del intervalo & Cola READY \\
\midrule
$0$--$2$   & P1 & $t=0.5$: arriba P2; $t=2$: P1 fin de quantum (rem.\ 6) & [P2(3), P1(6)] \\
$2$--$4$   & P2 & $t=3$: arriba P3; $t=4$: P2 fin de quantum (rem.\ 1)   & [P1(6), P3(5), P2(1)] \\
$4$--$6$   & P1 & $t=6$: P1 fin de quantum (rem.\ 4)                    & [P3(5), P2(1), P1(4)] \\
$6$--$8$   & P3 & $t=8$: P3 fin de quantum (rem.\ 3)                    & [P2(1), P1(4), P3(3)] \\
$8$--$9$   & P2 & $t=9$: P2 termina (ráfaga era 1)                      & [P1(4), P3(3)] \\
$9$--$11$  & P1 & $t=11$: P1 fin de quantum (rem.\ 2)                   & [P3(3), P1(2)] \\
$11$--$13$ & P3 & $t=13$: P3 fin de quantum (rem.\ 1)                   & [P1(2), P3(1)] \\
$13$--$15$ & P1 & $t=15$: P1 termina                                    & [P3(1)] \\
$15$--$16$ & P3 & $t=16$: P3 termina                                    & [\,] \\
\bottomrule
\end{tabular}
\end{center}

\textbf{Verificación}: la CPU está ocupada de 0 a 16 sin huecos y
$8 + 3 + 5 = 16$. Cuadra.

\subsection*{Diagrama de Gantt}

\begin{center}
\begin{tikzpicture}[xscale=0.45,yscale=0.7]
    \draw[fill=gp1]  (0,0)  rectangle (2,1)  node[midway,white,font=\bfseries\small]{P1};
    \draw[fill=gp2]  (2,0)  rectangle (4,1)  node[midway,white,font=\bfseries\small]{P2};
    \draw[fill=gp1]  (4,0)  rectangle (6,1)  node[midway,white,font=\bfseries\small]{P1};
    \draw[fill=gp3]  (6,0)  rectangle (8,1)  node[midway,white,font=\bfseries\small]{P3};
    \draw[fill=gp2]  (8,0)  rectangle (9,1)  node[midway,white,font=\bfseries\small]{P2};
    \draw[fill=gp1]  (9,0)  rectangle (11,1) node[midway,white,font=\bfseries\small]{P1};
    \draw[fill=gp3]  (11,0) rectangle (13,1) node[midway,white,font=\bfseries\small]{P3};
    \draw[fill=gp1]  (13,0) rectangle (15,1) node[midway,white,font=\bfseries\small]{P1};
    \draw[fill=gp3]  (15,0) rectangle (16,1) node[midway,white,font=\bfseries\small]{P3};
    \foreach \x in {0,2,4,6,8,9,11,13,15,16}
        \draw (\x,0) -- (\x,-0.15) node[below,font=\scriptsize]{\x};
\end{tikzpicture}
\\\small\emph{Diagrama de Gantt --- Round Robin con $q = 2$}
\end{center}

\subsection{2.a Turnaround de cada proceso}

$T_i = \text{Fin}_i - \text{Arribo}_i$:

\begin{center}
\small
\begin{tabular}{@{}lccc@{}}
\toprule
Proceso & Arribo & Fin & Turnaround \\
\midrule
P1 & 0.0 & 15 & 15.0 \\
P2 & 0.5 & 9  & 8.5  \\
P3 & 3.0 & 16 & 13.0 \\
\bottomrule
\end{tabular}
\end{center}

\[
\overline{T}_{\text{RR}} = \frac{15 + 8.5 + 13}{3} = \frac{36.5}{3} \approx 12.17.
\]

\subsection{2.b Comparación con FCFS y SJF}

Los tres algoritmos no parten de las mismas hipótesis: FCFS y RR usan
los arribos originales; SJF (parte b del Ejercicio~1) asume todos en
$t = 0$. Aún así, la comparación es ilustrativa:

\begin{center}
\small
\begin{tabular}{@{}lccc@{}}
\toprule
Proceso & $T$ FCFS & $T$ SJF & $T$ RR \\
\midrule
P1 & 8.0  & 16  & 15.0 \\
P2 & 10.5 & 3   & 8.5  \\
P3 & 13.0 & 8   & 13.0 \\
\midrule
\textbf{Promedio} & \textbf{10.50} & \textbf{9.00} & \textbf{12.17} \\
\bottomrule
\end{tabular}
\end{center}

Lectura:

\begin{itemize}[leftmargin=*,nosep]
    \item \textbf{SJF} es el mejor en turnaround promedio (9.0): la
        propiedad de optimalidad garantiza que ejecutar primero las
        ráfagas más cortas minimiza la espera promedio.
    \item \textbf{FCFS} queda intermedio (10.5). Sufre el
        \emph{convoy effect} porque P1 llega primero y arrastra a los
        demás, pero cada proceso ejecuta de un solo tirón.
    \item \textbf{RR} es el peor en promedio (12.17): el quantum
        chico fragmenta la ejecución de P1 en cuatro pedazos, pero a
        cambio P2 mejora notablemente (10.5 $\to$ 8.5) porque no
        queda atascada detrás de la ráfaga más larga.
\end{itemize}

Esto es justamente el compromiso clásico que la teoría señala: RR no
está pensado para minimizar turnaround promedio sino para acotar el
\emph{tiempo de respuesta} en sistemas \emph{time-sharing}. Con
quantum tendiendo al infinito, RR converge a FCFS; con quantum muy
chico, el overhead de cambios de contexto domina.

% =====================================================================
\section{Ejercicio 3 --- Prioridades dinámicas (estilo 4.3BSD)}
% =====================================================================

Un planificador recalcula cada segundo la prioridad de cada proceso
con la fórmula
\[
\text{priority} = \frac{\text{recent\_cpu\_usage}()}{2} + 60
\]
donde \emph{a mayor valor, menor prioridad}. Tres procesos tienen
usos recientes de CPU de 30, 15 y 10 respectivamente.

\subsection*{Cálculo de prioridades}

\begin{center}
\small
\begin{tabular}{@{}lccl@{}}
\toprule
Proceso & recent\_cpu & priority $=$ recent\_cpu$/2$ $+$ 60 & Posición \\
\midrule
P1 & 30 & $30/2 + 60 = \mathbf{75}$    & menor prioridad \\
P2 & 15 & $15/2 + 60 = \mathbf{67.5}$  & intermedia \\
P3 & 10 & $10/2 + 60 = \mathbf{65}$    & mayor prioridad \\
\bottomrule
\end{tabular}
\end{center}

Orden de prioridad (de más alta a más baja): $P_3 \succ P_2 \succ P_1$.

\subsection{3.a CPU-bound vs.\ I/O-bound}

Un proceso \textbf{CPU-bound} ejecuta ráfagas largas: gasta su
quantum sin bloquearse, así que su contador
\texttt{recent\_cpu\_usage} crece rápido. Su \texttt{priority}
numérica sube y, como mayor valor implica menor prioridad efectiva,
\textbf{el algoritmo lo posterga}. P1 (uso 30, prioridad 75) es el
ejemplo en la tabla.

Un proceso \textbf{I/O-bound} se bloquea seguido en \texttt{read} o
\texttt{write} y rara vez consume el quantum entero. Su contador
queda chico, su \texttt{priority} se mantiene cerca del piso (60), y
el algoritmo \textbf{lo favorece} cuando vuelve de I/O. P3 (uso 10,
prioridad 65) lo ilustra.

El término ``$+60$'' funciona como piso fijo (ningún proceso baja de
60) y el ``$/2$'' amortigua el crecimiento, evitando que un pico
breve de uso castigue desproporcionadamente al proceso. La fórmula
provista es exactamente la receta del scheduler clásico de
\textbf{4.3BSD UNIX}: \emph{multilevel feedback} con \emph{decay}
exponencial del \texttt{recent\_cpu}.

\subsection{3.b Objetivo del planificador}

Dos objetivos complementarios:

\begin{itemize}[leftmargin=*,nosep]
    \item \textbf{Mejor responsiveness para procesos interactivos}:
        editores, shells y servicios atendiendo \emph{requests} son
        típicamente I/O-bound. Que el scheduler les dé prioridad alta
        cuando vuelven de un bloqueo significa baja latencia entre el
        evento y su procesamiento.
    \item \textbf{Mejor uso de los dispositivos de I/O}: si los
        procesos I/O-bound se planifican rápido mantienen los
        dispositivos ocupados; los CPU-bound usan la CPU mientras los
        otros están bloqueados, sin perjudicar el throughput global.
\end{itemize}

La penalización a los CPU-bound \textbf{no es permanente}: como
\texttt{recent\_cpu\_usage} solo refleja la historia reciente, un
proceso CPU-bound que deja de demandar CPU verá caer su prioridad
numérica y volverá a ser elegible. Es exactamente el rol que la
teoría le adjudica al \emph{aging}.

En una línea: el planificador busca un \textbf{compromiso entre
utilización de CPU, throughput y tiempo de respuesta}, que la teoría
identifica como el objetivo central de los \emph{short-term
schedulers} en sistemas \emph{time-sharing}.

% =====================================================================
\section{Ejercicio 4 --- Planificación por prioridades estáticas}
% =====================================================================

Para los procesos del Ejercicio~1 (todos arriban en $t = 0$), con
prioridades:

\begin{center}
\small
\begin{tabular}{@{}lcc@{}}
\toprule
Proceso & CPU & Prioridad \\
\midrule
P1 & 8 & 1 \\
P2 & 3 & 3 \\
P3 & 5 & 2 \\
\bottomrule
\end{tabular}
\end{center}

\paragraph{Convención.} El enunciado no fija si menor número implica
mayor o menor prioridad. Se adopta la convención más extendida en
bibliografía clásica (Silberschatz, Tanenbaum) y la que usa Linux en
el rango \emph{real-time}: \textbf{menor número $\to$ mayor
prioridad}. Bajo esta convención el orden decreciente es
$P_1 \succ P_3 \succ P_2$.

\paragraph{Preemptivo vs.\ no preemptivo.} Como los tres llegan en
$t = 0$ y nadie se incorpora después, no hay diferencia entre las dos
versiones: una vez decidido el orden, ningún arribo posterior lo
puede alterar.

\subsection*{Diagrama de Gantt}

\begin{center}
\begin{tikzpicture}[xscale=0.5,yscale=0.7]
    \draw[fill=gp1] (0,0)  rectangle (8,1)  node[midway,white,font=\bfseries]{P1};
    \draw[fill=gp3] (8,0)  rectangle (13,1) node[midway,white,font=\bfseries]{P3};
    \draw[fill=gp2] (13,0) rectangle (16,1) node[midway,white,font=\bfseries]{P2};
    \foreach \x in {0,8,13,16}
        \draw (\x,0) -- (\x,-0.15) node[below,font=\small]{\x};
\end{tikzpicture}
\\\small\emph{Diagrama de Gantt --- Prioridades estáticas (menor número $\to$ mayor prioridad)}
\end{center}

\begin{center}
\small
\begin{tabular}{@{}lcccc@{}}
\toprule
Proceso & Inicio & Fin & Espera & Turnaround \\
\midrule
P1 & 0  & 8  & 0  & 8  \\
P3 & 8  & 13 & 8  & 13 \\
P2 & 13 & 16 & 13 & 16 \\
\bottomrule
\end{tabular}
\end{center}

Promedios:
$\overline{W} = (0 + 8 + 13)/3 = 7$;
$\overline{T} = (8 + 13 + 16)/3 \approx 12{.}33$.

\subsection*{Observaciones}

\begin{enumerate}[leftmargin=*,nosep]
    \item La secuencia depende \emph{exclusivamente} del número de
        prioridad. Aquí P1 (mayor prioridad) coincide con la ráfaga
        más larga, lo que perjudica el promedio respecto de SJF
        (orden P2-P3-P1).
    \item La planificación por prioridades \textbf{no garantiza
        optimalidad} en tiempo de espera ni en turnaround promedio:
        solo respeta la importancia relativa que el sistema asigna a
        cada proceso. Es una decisión política, no temporal.
    \item \textbf{Riesgo de starvation}: si llegasen continuamente
        procesos con prioridad 1 o 2, P2 (prioridad 3) podría
        postergarse indefinidamente. La solución estándar
        ---mencionada explícitamente por la teoría--- es el
        \emph{aging}: incrementar gradualmente la prioridad de
        procesos que llevan mucho tiempo esperando.
\end{enumerate}

\subsection*{Si la convención fuera ``mayor número $\to$ mayor prioridad''}

El orden sería $P_2 \succ P_3 \succ P_1$:

\begin{center}
\begin{tikzpicture}[xscale=0.5,yscale=0.7]
    \draw[fill=gp2] (0,0) rectangle (3,1)   node[midway,white,font=\bfseries]{P2};
    \draw[fill=gp3] (3,0) rectangle (8,1)   node[midway,white,font=\bfseries]{P3};
    \draw[fill=gp1] (8,0) rectangle (16,1)  node[midway,white,font=\bfseries]{P1};
    \foreach \x in {0,3,8,16}
        \draw (\x,0) -- (\x,-0.15) node[below,font=\small]{\x};
\end{tikzpicture}
\\\small\emph{Diagrama de Gantt --- convención inversa}
\end{center}

Promedios: $\overline{W} \approx 3{.}67$,
$\overline{T} = 9$. Coincide numéricamente con SJF en el
Ejercicio~1, porque casualmente las ``prioridades altas'' coinciden
con las ráfagas más cortas.

% =====================================================================
\section{Ejercicio 5 --- Multilevel Feedback de tres niveles}
% =====================================================================

Se pide describir paso a paso un scheduler MLF con tres niveles
$L_0$, $L_1$, $L_2$ y \emph{quantums} 3, 2 y 1 ms respectivamente,
sobre las tareas:

\begin{center}
\small
\begin{tabular}{@{}lp{8cm}@{}}
\toprule
Tarea & Comportamiento (línea temporal) \\
\midrule
T1 & CPU $=$ 5 ms ; WAIT $=$ 5 ms ; CPU $=$ 1 ms \\
T2 & CPU $=$ 12 ms \\
T3 & CPU $=$ 1 ms ; WAIT $=$ 2 ms ; CPU $=$ 1 ms \\
\bottomrule
\end{tabular}
\end{center}

Las tres ingresan a $L_0$ en $t = 0$ en el orden T1, T2, T3.

\subsection*{Reglas del MLF}

\begin{itemize}[leftmargin=*,nosep]
    \item Niveles ordenados por prioridad: $L_0$ es el más
        prioritario, $L_2$ el menos.
    \item El scheduler ejecuta siempre al primero del nivel más alto
        \emph{no vacío}.
    \item Dentro de cada nivel se aplica RR con su propio quantum.
    \item \textbf{Demote}: si un proceso \emph{consume todo el
        quantum} se considera CPU-bound y baja un nivel. Si ya está
        en $L_2$ se reencola en el mismo $L_2$.
    \item \textbf{Sin demote}: si un proceso \emph{cede
        voluntariamente} (bloqueo en I/O) o \emph{termina} antes del
        quantum, se considera interactivo y se queda en su nivel.
    \item \textbf{Preemption entre niveles}: si mientras un proceso
        ejecuta llega/despierta otro de nivel superior, el primero es
        desalojado y reencolado al final de su nivel actual.
    \item Convención en empates: (a) si llegada y desalojo coinciden
        en el tiempo, primero se encola la nueva; (b) si la ráfaga
        termina justo al expirar el quantum, prevalece la lectura
        ``I/O'' (no demote); (c) un proceso que despierta de WAIT
        regresa al mismo nivel donde estaba antes del bloqueo.
\end{itemize}

\subsection*{Tabla temporal}

Notación: \texttt{Ti:n} significa ``tarea $i$ con $n$ ms pendientes
de CPU en su ráfaga actual''. \texttt{<>} es cola vacía.

\begin{center}
\scriptsize
\begin{tabular}{@{}llllllllp{1.6cm}@{}}
\toprule
Time & CPU & T1 & T2 & T3 & $L_0$ & $L_1$ & $L_2$ \\
\midrule
0     & ---             & RDY  & RDY & RDY  & \texttt{<T1:5,T2:12,T3:1>} & \texttt{<>}        & \texttt{<>}      \\
0--3  & T1              & RUN  & RDY & RDY  & \texttt{<T2:12,T3:1>}      & \texttt{<>}        & \texttt{<>}      \\
3     & demote T1       & RDY  & RDY & RDY  & \texttt{<T2:12,T3:1>}      & \texttt{<T1:2>}    & \texttt{<>}      \\
3--6  & T2              & RDY  & RUN & RDY  & \texttt{<T3:1>}            & \texttt{<T1:2>}    & \texttt{<>}      \\
6     & demote T2       & RDY  & RDY & RDY  & \texttt{<T3:1>}            & \texttt{<T1:2,T2:9>} & \texttt{<>}    \\
6--7  & T3              & RDY  & RDY & RUN  & \texttt{<>}                & \texttt{<T1:2,T2:9>} & \texttt{<>}    \\
7     & T3 $\to$ WAIT   & RDY  & RDY & WAIT & \texttt{<>}                & \texttt{<T1:2,T2:9>} & \texttt{<>}    \\
7--9  & T1              & RUN  & RDY & WAIT & \texttt{<>}                & \texttt{<T2:9>}    & \texttt{<>}      \\
9     & T1 $\to$ WAIT; T3 wakeup & WAIT & RDY & RDY & \texttt{<T3:1>}    & \texttt{<T2:9>}    & \texttt{<>}      \\
9--10 & T3              & WAIT & RDY & RUN  & \texttt{<>}                & \texttt{<T2:9>}    & \texttt{<>}      \\
10    & T3 termina      & WAIT & RDY & TERM & \texttt{<>}                & \texttt{<T2:9>}    & \texttt{<>}      \\
10--12 & T2             & WAIT & RUN & TERM & \texttt{<>}                & \texttt{<>}        & \texttt{<>}      \\
12    & demote T2       & WAIT & RDY & TERM & \texttt{<>}                & \texttt{<>}        & \texttt{<T2:7>}  \\
12--13 & T2             & WAIT & RUN & TERM & \texttt{<>}                & \texttt{<>}        & \texttt{<>}      \\
13    & T2 reencola en $L_2$ & WAIT & RDY & TERM & \texttt{<>}           & \texttt{<>}        & \texttt{<T2:6>}  \\
13--14 & T2             & WAIT & RUN & TERM & \texttt{<>}                & \texttt{<>}        & \texttt{<>}      \\
14    & T1 wakeup $\to L_1$; T2 desalojado & RDY & RDY & TERM & \texttt{<>} & \texttt{<T1:1>} & \texttt{<T2:5>} \\
14--15 & T1             & RUN  & RDY & TERM & \texttt{<>}                & \texttt{<>}        & \texttt{<T2:5>}  \\
15    & T1 termina      & TERM & RDY & TERM & \texttt{<>}                & \texttt{<>}        & \texttt{<T2:5>}  \\
15--20 & T2             & TERM & RUN & TERM & \texttt{<>}                & \texttt{<>}        & \texttt{<T2:4..0>} \\
20    & T2 termina      & TERM & TERM & TERM & \texttt{<>}               & \texttt{<>}        & \texttt{<>}      \\
\bottomrule
\end{tabular}
\end{center}

\textbf{Verificación}: la CPU está ocupada de 0 a 20 ms sin huecos.
Las ráfagas suman $T_1 = 6$, $T_2 = 12$, $T_3 = 2$, total 20 ms.

\subsection*{Justificación de los pasos clave}

\begin{itemize}[leftmargin=*,nosep]
    \item \textbf{$t = 0$--$3$}: T1 toma CPU en $L_0$ (quantum 3).
        Su ráfaga era 5; consume el quantum, le quedan 2. Demote
        a $L_1$.
    \item \textbf{$t = 3$--$6$}: T2 corre los 3 ms del quantum de
        $L_0$. Le quedan 9. Demote a $L_1$.
    \item \textbf{$t = 6$--$7$}: T3 obtiene CPU en $L_0$ y completa
        su ráfaga de 1 ms (cede voluntariamente al pasar a WAIT por 2
        ms). \textbf{No demote}: cuando despierte vuelve a $L_0$.
    \item \textbf{$t = 7$--$9$}: con $L_0$ vacío, el scheduler baja a
        $L_1$ y elige T1. T1 corre los 2 ms del quantum de $L_1$ y
        \emph{justo} completa su primera ráfaga de 5 ms (3 en $L_0$
        + 2 en $L_1$); pasa a WAIT por 5 ms. Por la convención
        adoptada (ráfaga termina exactamente en el quantum $\to$
        I/O), T1 no demote a $L_2$: cuando despierte volverá a $L_1$.
        Simultáneamente, T3 despierta y vuelve a $L_0$.
    \item \textbf{$t = 9$--$10$}: T3 obtiene CPU (preferencia de
        $L_0$ sobre $L_1$), corre su última ráfaga de 1 ms y termina.
    \item \textbf{$t = 10$--$13$}: T2 ejecuta en $L_1$ (2 ms), demote
        a $L_2$, ejecuta otro ms en $L_2$. En cada quantum agotado
        en $L_2$ se reencola en el mismo nivel.
    \item \textbf{$t = 13$--$14$}: T2 sigue en $L_2$. Justo al
        expirar este quantum, T1 despierta y entra a $L_1$.
    \item \textbf{$t = 14$--$15$}: T1 obtiene CPU (preferencia de
        $L_1$ sobre $L_2$), corre 1 ms y termina su última ráfaga.
    \item \textbf{$t = 15$--$20$}: T2 queda sola y completa los 5 ms
        restantes en $L_2$ (cinco quantums consecutivos de 1 ms).
\end{itemize}

\subsection*{Métricas resultantes}

\begin{center}
\small
\begin{tabular}{@{}lccccc@{}}
\toprule
Tarea & Llegada & Fin & Turnaround & CPU usada & Espera \\
\midrule
T1 & 0 & 15 & 15 & 6  & 4 \\
T2 & 0 & 20 & 20 & 12 & 8 \\
T3 & 0 & 10 & 10 & 2  & 6 \\
\bottomrule
\end{tabular}
\end{center}

\textbf{Lectura.} T3 (interactivo, 2 ms de CPU) sale rápido del
sistema; T1 (mixta, 6 ms de CPU) termina segunda. T2 (CPU-bound puro,
12 ms) es \emph{la más castigada}: queda atrapada en $L_2$ y solo
termina al final, después de que las otras se hayan ido. Sin embargo
\textbf{no padece starvation} porque, una vez que el sistema queda
libre, T2 obtiene la CPU sin contención. Es exactamente el
comportamiento que la teoría atribuye al MLF: \emph{mejorar el tiempo
de respuesta de los procesos interactivos castigando a los CPU-bound,
sin negarles servicio}.

% =====================================================================
\section{Ejercicio 6 --- Processor affinity vs.\ load balancing en SMP}
% =====================================================================

En sistemas SMP el scheduler debe lograr \emph{processor affinity} y
\emph{load balancing} simultáneamente. Se pide explicar por qué estas
propiedades se contraponen.

\subsection*{Definiciones}

\begin{itemize}[leftmargin=*,nosep]
    \item \textbf{Load balancing}: emparejar la carga de trabajo
        entre todas las CPUs. En cada \emph{context switch} un
        proceso puede \emph{migrar} de CPU.
    \item \textbf{Processor affinity}: minimizar la migración de
        procesos entre CPUs. Linux ofrece la syscall
        \texttt{sched\_setaffinity} (\emph{hard affinity}) para
        ligar un proceso a un subconjunto de CPUs.
\end{itemize}

La teoría explicita el motivo: la migración \emph{afecta al
rendimiento porque la memoria caché de la CPU anterior queda
invalidada y la caché de la nueva debe volver a llenarse}.

\subsection*{Por qué se contraponen}

Las dos propiedades empujan al scheduler en direcciones opuestas:

\begin{enumerate}[leftmargin=*,nosep]
    \item \textbf{Load balancing pide migrar.} Si una CPU está
        sobrecargada (cola READY larga) y otra está libre, lo correcto
        \emph{desde el balance} es mover algún proceso a la libre.
        Sin migración, una CPU queda al 100\% mientras otra
        desperdicia ciclos en idle.
    \item \textbf{Processor affinity pide no migrar.} Cada migración
        descarta el trabajo de \emph{calentamiento de caché} (L1 y L2
        son privadas por core). El proceso paga miles de instrucciones
        con \emph{cache miss} hasta repoblar la nueva caché. También
        se pierde la huella en la TLB.
\end{enumerate}

El conflicto concreto:

\begin{itemize}[leftmargin=*,nosep]
    \item Si el scheduler \textbf{prioriza affinity}, deja de migrar
        incluso con desbalance: aparece subutilización (una CPU al
        100\%, otra al 0\%). El sistema de $N$ CPUs rinde como uno
        más chico.
    \item Si el scheduler \textbf{prioriza load balancing}, migra
        seguido para mantener parejas las CPUs. Pero cada migración
        paga el costo de caché. Si las migraciones son frecuentes el
        sistema gasta tiempo en repoblar caches en lugar de hacer
        trabajo útil.
\end{itemize}

\subsection*{Compromiso práctico}

La teoría lo nombra explícitamente: ``estos dos objetivos se
contraponen, por lo que hay que lograr una solución de
compromiso''. Las estrategias clásicas son:

\begin{itemize}[leftmargin=*,nosep]
    \item \textbf{Per-CPU runqueue} más un kernel thread dedicado al
        balanceo (en Linux: el thread \texttt{migration} y la rutina
        \texttt{load\_balance}). El scheduler local respeta la
        affinity por defecto; el balanceador estima la carga y
        migra solo si el desbalance excede un umbral.
    \item \textbf{Soft affinity}: el scheduler trata de no migrar
        pero puede hacerlo si compensa (default en Linux).
    \item \textbf{Hard affinity}: con \texttt{sched\_setaffinity} el
        proceso queda restringido a un conjunto fijo de CPUs.
    \item \textbf{Push vs.\ pull migration}: una CPU sobrecargada
        empuja trabajo, o una CPU ociosa lo tira. Coexisten.
\end{itemize}

El principio común: \textbf{migrar solo cuando el desbalance
proyectado supera al costo esperado de la migración}. En el extremo
NUMA, donde la migración entre nodos suma además el costo de
accesos a memoria remota, la affinity gana peso. En arquitecturas
\emph{multi-core} con caché L3 compartida entre cores del mismo
\emph{socket}, migrar dentro del socket es mucho más barato que
entre sockets, y los schedulers modernos (CFS de Linux) son
conscientes de esta jerarquía.

\paragraph{Resumen.} \emph{Affinity} protege la caché de cada
proceso individual; \emph{load balancing} protege la utilización
global del multiprocesador. Cada migración gana balance al costo de
afinidad y viceversa, así que el scheduler decide cuánto vale cada
migración antes de hacerla.

% =====================================================================
\section{Conclusiones}
% =====================================================================

El práctico recorre el menú clásico de algoritmos de planificación de
CPU y deja en evidencia los compromisos centrales del diseño:

\begin{enumerate}[leftmargin=*,nosep]
    \item \textbf{FCFS} es simple y predecible pero sufre el
        \emph{convoy effect} cuando una ráfaga larga llega primero
        (Ej.~1).
    \item \textbf{SJF} minimiza el tiempo de espera promedio cuando
        se conocen las ráfagas (Ej.~1), pero es no preemptivo y
        requiere conocer las duraciones (estimable con promedio
        exponencial $\tau_{n+1} = \alpha t_n + (1 - \alpha)\tau_n$).
    \item \textbf{Round Robin} acota el tiempo de respuesta a costa
        de un mayor turnaround promedio cuando hay procesos largos
        (Ej.~2). El quantum es el parámetro de compromiso.
    \item \textbf{Prioridades} permiten al sistema imponer una
        política, pero no garantizan optimalidad temporal y traen el
        riesgo de \emph{starvation} (Ej.~3 y 4). El \emph{aging} es
        la solución estándar.
    \item \textbf{Multilevel Feedback} combina prioridades dinámicas
        con RR por nivel; es la versión ``natural'' para sistemas
        \emph{time-sharing} con cargas heterogéneas, y favorece a los
        procesos interactivos sin negarle progreso a los CPU-bound
        (Ej.~5).
    \item En \textbf{SMP} el scheduler debe equilibrar affinity y
        load balancing: dos objetivos legítimos pero opuestos, que
        obligan a tomar decisiones costo-beneficio en cada migración
        (Ej.~6).
\end{enumerate}

La idea unificadora es que la planificación es siempre un
compromiso entre criterios que se contraponen: utilización,
\emph{throughput}, tiempo de espera, tiempo de respuesta y
\emph{fairness}. Ningún algoritmo gana en todas las métricas a la
vez; cada uno tiene su nicho de aplicación. Los schedulers modernos
(CFS, EEVDF en Linux desde 2024) son combinaciones sofisticadas de
estos principios, sintonizados para cargas reales heterogéneas.

\end{document}
